\section{Discussion}
\label{sec:discussion}

The experiments distinguish predictive performance from governance performance. HOG--SVM reduced cell error to 0.19\% in the synthetic evaluation, but five mistakes remained. The authority boundary addresses the destination of such errors: even a selected prediction is persisted as a candidate until an attributable, version-checked review creates a fact. Stronger perception therefore reduces reviewer workload without making candidate--fact separation redundant.

Recognizer choice also changed the role of selective prediction. The template matcher needed to route 88.81\% of evaluation rows to achieve zero observed accepted errors, whereas HOG--SVM satisfied the calibration risk target at full coverage. This difference argues against treating a threshold as a fixed safety setting. Selection must be recalibrated after changes in forms, printers, cameras, or recognizers, and a deployment threshold should reflect the relative cost of an unreviewed error, reviewer time, and delayed processing.

The whole-form benchmark exposed failures that the cell benchmark could not. Digit and OMR recognition were perfect under blurred forms once fields were localized, yet QR identification failed, so no complete record could be produced. Under rotation, ArUco alignment failed and digit accuracy fell sharply. These observations make the engineering priorities concrete: QR preprocessing or fallback classification is needed for blur, and marker detection or orientation normalization is needed for rotation. The 60\% complete-form result is consequently more informative than the 99.81\% isolated-cell score for the current pipeline.

The controlled fault cases, candidate-isolation ablation, and 500 randomized trials test the causal role of capability separation. The incorrect value was harmless when stored through a candidate-only interface and changed the fact when the experiment deliberately exposed the review capability. This is stronger evidence than an architecture diagram alone, but weaker than formal noninterference verification, database privilege separation, or adversarial assessment of every UI and migration path.

The authority-integration benchmark tests the boundary in its executable form. Certificates bind a candidate to its field, evidence, and expected version, so a certificate moved to another field or confirmed against another form's evidence is rejected before any write. Confirmations commit per-field decisions, transitions, and the record version in one compare-and-swap transaction, so concurrent and stale confirmations fail closed. Reverse traces are rebuilt from the database per field rather than assumed, and the benchmark shows that tampered producers, deleted transitions, and record-version rebinding---including rebinding to another existing row that the database foreign key still accepts---are reported as incomplete. The measured silent fault escape of 0.75 belongs to the scripted reviewer and is deliberately not interpreted as human error rates; that question is deferred to the controlled human-review study.


The architecture remains model-agnostic. A conformal or learned selector can replace the decision-margin rule, and different retrieval or language models can supply assistance, provided that each terminates outside the fact-writing interface. This separation allows model procurement to change while preserving a stable audit question: which service can invoke the fact-transition operation?

The next empirical study should use a blinded, double-entered corpus of real forms with adjudicated ground truth and subgroup reporting by template, device, writing style, and degradation. A second study should randomize review interfaces---recognition only, retrieval-assisted, and AI-assisted---and measure review time, correction accuracy, automation bias, and inappropriate suggestion acceptance. Those studies would test decision quality and operational value, which synthetic execution cannot establish.
